\newif\ifglobal

%%%%%%%%%%%%%%%%%%%%%%%
%%% Compile to view %%%
%%%%%%%%%%%%%%%%%%%%%%%

%%%%%%%%%%%%%%%%%%%%%%%%%%%%%%%%%%%%%%%%%%%%%%%%%%%
%%% Readme:                                     %%%
%%% https://www.overleaf.com/read/bwdjvfjksvjy/ %%%
%%%%%%%%%%%%%%%%%%%%%%%%%%%%%%%%%%%%%%%%%%%%%%%%%%%

% >>> M?: marks a module setting number or important notice

% >>> M4: Set {./relative-path} to external files for this module <<<
\def \modulefiles {./18-SYSREG}
% To add an external file, use:
% (Replace '---' with actual filename)
% '\input{\modulefiles/tables/---.tex}' for tables
% '\includegraphics{\modulefiles/figures/---}' for figures
% '\subfile{\modulefiles/---}' for register files


% >>> M1: This module is in global project? <<<
% 'YES' - uncomment; 'NO' - comment out
\globaltrue

\ifglobal
    \documentclass[main\_\_EN.tex]{subfiles}

\else
    \documentclass[openany, 10pt]{book}

    % >>> M2: In ./00-shared/preamble-trm-module, also find
    %         settings for visibility of todo notes, watermark <<<
    \usepackage{preamble-trm-module}

    % >>> M3: Temporary labels in English,
    %         you can add your own labels there <<<
    \myexternaldocument{./00-shared/temp-labels-en}

\fi %\ifglobal


\setcounter{secnumdepth}{4}


% >>> M5: If this module needs any updates in
% './00-shared/preamble-trm-module.sty'
% (include additional package, etc.), contact your module PM
% or create the file './00-shared/changelog.tex' make the
% necessary updates yourself and document them in this file <<<


\begin{document}


% Sets language into which pre-defined bits of text are translated
\selectlanguage{english}

% Do not delete this block
\ifglobal
% Add the following front matter if
% this module is outside of global project
\else
    \InputIfFileExists{readme}{}{}
    {\let\clearpage\relax \listoftodos}
    \tableofcontents
    %\listoftables
    %\listoffigures
\fi


%%%%%%%%%%%%%%%%%%%%%%%%%
%%% TRM Module Begins %%%
%%%%%%%%%%%%%%%%%%%%%%%%%

\chapter{System Registers (SYSREG)}
\label{mod:sysreg}
\hypertarget{sysreg}{}

\section{Overview}

The \chipname{} integrates a large number of peripherals, and enables the control of individual peripherals to achieve optimal characteristics in performance-vs-power-consumption scenarios. Specifically, \chipname{} has various system configuration registers that can be used for the chip's clock management (clock gating), power management, and the configuration of peripherals and core-system modules. This chapter lists all these system registers and their functions.

\section{Features}

\chipname{} system registers can be used to control the following peripheral blocks and core modules:

\begin{itemize}
    \item System and memory
    \item Clock
    \item Software Interrupt
    \item Low-power management
    \item Peripheral clock gating and reset
\end{itemize}


\section{Function Description}

\subsection{System and Memory Registers}
\subsubsection{Internal Memory}

The following registers can be used to control \chipname{}'s internal memory:

\begin{itemize}
    \item In register \hyperref[regdesc:SYSCONCLKGATEFORCEONREG]{SYSCON\_CLKGATE\_FORCE\_ON\_REG}:
    \begin{itemize}
        \item Setting different bits of the \hyperref[fielddesc:SYSCONROMCLKGATEFORCEON]{SYSCON\_ROM\_CLKGATE\_FORCE\_ON} field forces on the clock gates of different blocks of Internal ROM 0 and Internal ROM 1.
        \item Setting different bits of the \hyperref[fielddesc:SYSCONSRAMCLKGATEFORCEON]{SYSCON\_SRAM\_CLKGATE\_FORCE\_ON} field forces on the clock gates of different blocks of Internal SRAM.
        \item This means when the respective bits of this register are set to 1, the clock gate of the corresponding ROM or SRAM blocks will always be on. Otherwise, the clock gate will turn on automatically when the corresponding ROM or SRAM blocks are accessed and turn off automatically when the corresponding ROM or SRAM blocks are not accessed. Therefore, it's recommended to configure these bits to 0 to lower power consumption.
    \end{itemize}
    \item In register \hyperref[regdesc:SYSCONMEMPOWERDOWNREG]{SYSCON\_MEM\_POWER\_DOWN\_REG}:
    \begin{itemize}
        \item Setting different bits of the \hyperref[fielddesc:SYSCONROMPOWERDOWN]{SYSCON\_ROM\_POWER\_DOWN} field sends different blocks of Internal ROM 0 and Internal ROM 1 into retention state.
        \item Setting different bits of the \hyperref[fielddesc:SYSCONSRAMPOWERDOWN]{SYSCON\_SRAM\_POWER\_DOWN} field sends different blocks of Internal SRAM into retention state.
        \item The ``Retention" state is a low-power state of a memory block. In this state, the memory block still holds all the data stored but cannot be accessed, thus reducing the power consumption. Therefore, you can send a certain block of memory into the retention state to reduce power consumption if you know you are not going to use such memory block for some time.
    \end{itemize}
    \item In register \hyperref[regdesc:SYSCONMEMPOWERUPREG]{SYSCON\_MEM\_POWER\_UP\_REG}:
    \begin{itemize}
        \item By default, all memory enters low-power state when the chip enters the Light-sleep mode.
        \item Setting different bits of the \hyperref[fielddesc:SYSCONROMPOWERUP]{SYSCON\_ROM\_POWER\_UP} field forces different blocks of Internal ROM 0 and Internal ROM 1 to work as normal (do not enter the retention state) when the chip enters Light-sleep.
        \item Setting different bits of the \hyperref[fielddesc:SYSCONSRAMPOWERUP]{SYSCON\_SRAM\_POWER\_UP} field forces different blocks of Internal SRAM to work as normal (do not enter the retention state) when the chip enters Light-sleep.
    \end{itemize}
\end{itemize}

For detailed information about the controlling bits of different blocks, please see Table \ref{tab:sysreg-mem-controlling-bit} below.

\begin{longtable}[c]{| l | l | l |l | l | l |}
\caption{Memory Controlling Bit}
\label{tab:sysreg-mem-controlling-bit}
\\\hline
\rowcolor{lightgray}
        Memory & Lowest Address1 & Highest Address1 & Lowest Address2 & Highest Address2 & Controlling Bit \\ \hline
        ROM 0  & 0x4000\_0000    & 0x4003\_FFFF  &- & -  & Bit0 \\\hline
        ROM 1  & 0x4004\_0000    & 0x4005\_FFFF  & 0x3FF0\_0000  & 0x3FF1\_FFFF & Bit1\\\hline
        SRAM Block 0  & 0x4037\_C000    & 0x4037\_FFFF  & - & -  & Bit0 \\\hline
        SRAM Block 1  & 0x4038\_0000    & 0x4039\_FFFF  & 0x3FC8\_0000  & 0x3FC9\_FFFF & Bit1\\\hline
        SRAM Block 2  & 0x403A\_0000    & 0x403B\_FFFF  & 0x3FCA\_0000  & 0x3FCB\_FFFF & Bit2\\\hline
        SRAM Block 3  & 0x403C\_0000    & 0x403D\_FFFF  & 0x3FCC\_0000  & 0x3FCD\_FFFF & Bit3\\\hline
\end{longtable}

For more information, please refer to Chapter \ref{mod:sysmem} \textit{\nameref{mod:sysmem}}.

\subsubsection{External Memory}

\hyperref[regdesc:SYSTEMEXTERNALDEVICEENCRYPTDECRYPTCONTROLREG]{SYSTEM\_EXTERNAL\_DEVICE\_ENCRYPT\_DECRYPT\_CONTROL\_REG} configures encryption and decryption options of the external memory. For details, please refer to Chapter \ref{mod:extmemencr}
\textit{\nameref{mod:extmemencr}}.

\subsubsection{RSA Memory}

\hyperref[regdesc:SYSTEMRSAPDCTRLREG]{SYSTEM\_RSA\_PD\_CTRL\_REG} controls the SRAM memory in the RSA accelerator.

\begin{itemize}
    \item Setting the \hyperref[fielddesc:SYSTEMRSAMEMPD]{SYSTEM\_RSA\_MEM\_PD} bit to send the RSA memory into retention state. This bit has the lowest priority, meaning it can be masked by the \hyperref[fielddesc:SYSTEMRSAMEMFORCEPU]{SYSTEM\_RSA\_MEM\_FORCE\_PU} field. This bit is invalid when the \textit{\nameref{mod:digsig}} occupies the RSA.
    \item Setting the \hyperref[fielddesc:SYSTEMRSAMEMFORCEPU]{SYSTEM\_RSA\_MEM\_FORCE\_PU} bit to force the RSA memory to work as normal when the chip enters light sleep. This bit has the second highest priority, meaning it overrides the \hyperref[fielddesc:SYSTEMRSAMEMPD]{SYSTEM\_RSA\_MEM\_PD} field.
    \item Setting the \hyperref[fielddesc:SYSTEMRSAMEMFORCEPD]{SYSTEM\_RSA\_MEM\_FORCE\_PD} bit to send the RSA memory into retention state. This bit has the highest priority, meaning it sends the RSA memory into retention state regardless of the \hyperref[fielddesc:SYSTEMRSAMEMFORCEPU]{SYSTEM\_RSA\_MEM\_FORCE\_PU} field.
\end{itemize}

\subsection{Clock Registers}

The following registers are used to set clock sources and frequency. For more information, please refer to Chapter \ref{mod:resclk}
\textit{\nameref{mod:resclk}}.

\begin{itemize}
    \item \hyperref[regdesc:SYSTEMCPUPERCONFREG]{SYSTEM\_CPU\_PER\_CONF\_REG}
    \item \hyperref[regdesc:SYSTEMSYSCLKCONFREG]{SYSTEM\_SYSCLK\_CONF\_REG}
    \item \hyperref[regdesc:SYSTEMBTLPCKDIVFRACREG]{SYSTEM\_BT\_LPCK\_DIV\_FRAC\_REG}
\end{itemize}

\subsection{Interrupt Signal Registers}

The following registers are used for generating the interrupt signals (software interrupt), which then can be routed to the CPU peripheral interrupts via the interrupt matrix. To be more specific, writing 1 to any of the following registers generates an interrupt signal. Therefore, these registers can be used by software to control interrupts. The following registers correspond to the interrupt source SW\_INTR\_0/1/2/3. For more information, please refer to Chapter \ref{mod:intmtrx}
\textit{\nameref{mod:intmtrx}}.

\begin{itemize}
    \item \hyperref[regdesc:SYSTEMCPUINTRFROMCPU0REG]{SYSTEM\_CPU\_INTR\_FROM\_CPU\_0\_REG}
    \item \hyperref[regdesc:SYSTEMCPUINTRFROMCPU1REG]{SYSTEM\_CPU\_INTR\_FROM\_CPU\_1\_REG}
    \item \hyperref[regdesc:SYSTEMCPUINTRFROMCPU2REG]{SYSTEM\_CPU\_INTR\_FROM\_CPU\_2\_REG}
    \item \hyperref[regdesc:SYSTEMCPUINTRFROMCPU3REG]{SYSTEM\_CPU\_INTR\_FROM\_CPU\_3\_REG}
\end{itemize}

\subsection{Low-power Management Registers}

The following registers are used for low-power management. For more information, please refer to Chapter \ref{mod:lowpowm} \textit{\nameref{mod:lowpowm}}.
\begin{itemize}
    \item \hyperref[regdesc:SYSTEMRTCFASTMEMCONFIGREG]{SYSTEM\_RTC\_FASTMEM\_CONFIG\_REG}: configures the RTC CRC check.
    \item \hyperref[regdesc:SYSTEMRTCFASTMEMCRCREG]{SYSTEM\_RTC\_FASTMEM\_CRC\_REG}: configures the CRC check value.
\end{itemize}

\subsection{Peripheral Clock Gating and Reset Registers}

The following registers are used for controlling the clock gating and reset of different peripherals. Details can be seen in Table \ref{tab:sysreg-peripheral-controlling-bit}.

\begin{itemize}
    \item \hyperref[regdesc:SYSTEMCACHECONTROLREG]{SYSTEM\_CACHE\_CONTROL\_REG}
    \item \hyperref[regdesc:SYSTEMPERIPCLKEN0REG]{SYSTEM\_PERIP\_CLK\_EN0\_REG}
    \item \hyperref[regdesc:SYSTEMPERIPRSTEN0REG]{SYSTEM\_PERIP\_RST\_EN0\_REG}
    \item \hyperref[regdesc:SYSTEMPERIPCLKEN1REG]{SYSTEM\_PERIP\_CLK\_EN1\_REG}
    \item \hyperref[regdesc:SYSTEMPERIPRSTEN1REG]{SYSTEM\_PERIP\_RST\_EN1\_REG}

\end{itemize}

\chipname{} features low power consumption. This is why some peripheral clocks are gated (disabled) by default. Before using any of these peripherals, it is mandatory to enable the clock for the given peripheral and release the peripheral from reset state. For details, see the table below:

\begin{ThreePartTable}

\begin{TableNotes}
    \item[1] Set the clock enable bit to 1 to enable the clock, and to 0 to disable the clock;
    \item[2] Set the reset enabling bit to 1 to reset a peripheral, and to 0 to disable the reset.
    \item[3] Reset registers cannot be cleared by hardware. Therefore, SW reset clear is required after setting the reset registers.
    \item[4] UART memory is shared by all UART peripherals, meaning having any active UART peripherals will prevent the UART memory from entering the clock-gated state.
    \item[5] When DMA is required for periphral communications, for example, UCHI0, SPI2, I2S, AES, SHA, and ADC, DMA clock should also be enabled.
    \item[6] Resetting this bit also resets the SHA accelerator.
    \item[7] Resetting this bit also resets the AES, SHA, and RSA accelerators.

\end{TableNotes}

\begin{longtable}[c]{ |l|l|l| }
    \caption{Clock Gating and Reset Bits}
    \label{tab:sysreg-peripheral-controlling-bit}
    \\\hline\rowcolor{lightgray}
    \textbf{Component}    & \textbf{Clock Enabling Bit \tnote{1}} & \textbf{Reset Controlling Bit \tnote{2} \tnote{3}}\\\hline
    \endfirsthead
    \multicolumn{3}{c}{\bfseries \tablename\ \thetable{} -- cont'd from previous page}\\\hline
    \rowcolor{lightgray}
    \textbf{Component}    & \textbf{Clock Enabling Bit \tnote{1}} & \textbf{Reset Controlling Bit \tnote{2} \tnote{3}}\\\hline
    \endhead
    \multicolumn{3}{r}{\textbf{Cont'd on next page}} \\
    \endfoot
    \insertTableNotes
    \endlastfoot
%%%%%%%%%%%%%%%%%%%%%%%%%%%%%%%%%%%%%%%%%%%%%%%%%%%%%%%%%%%%%%%%%%%%%%%%%%%%%%%%%%%%%%%%%%%%%%%%%%%%%%%%%%%%%%%%%%%%%%%%%%%

%%%%%%%%%%%%%%%%%%%%%%%%%%%%%%%%%%%%%%%%%%%%%%%%%%%%%%%%%%%%%%%%%%%%%%%%%%%%%%%%%%%%%%%%%%%%%%%%%%%%%%%%%%%%%%%%%%%%%%%%%%%
    \textbf{CACHE Control} & \multicolumn{2}{|c|}{\textbf{\hyperref[regdesc:SYSTEMCACHECONTROLREG]{SYSTEM\_CACHE\_CONTROL\_REG}}} \\\hline

    DCACHE & \hyperref[fielddesc:SYSTEMDCACHECLKON]{SYSTEM\_DCACHE\_CLK\_ON} & \hyperref[fielddesc:SYSTEMDCACHERESET]{SYSTEM\_DCACHE\_RESET} \\\hline
    ICACHE & \hyperref[fielddesc:SYSTEMICACHECLKON]{SYSTEM\_ICACHE\_CLK\_ON} & \hyperref[fielddesc:SYSTEMICACHERESET]{SYSTEM\_ICACHE\_RESET} \\\hline

%%%%%%%%%%%%%%%%%%%%%%%%%%%%%%%%%%%%%%%%%%%%%%%%%%%%%%%%%%%%%%%%%%%%%%%%%%%%%%%%%%%%%%%%%%%%%%%%%%%%%%%%%%%%%%%%%%%%%%%%%%%

%%%%%%%%%%%%%%%%%%%%%%%%%%%%%%%%%%%%%%%%%%%%%%%%%%%%%%%%%%%%%%%%%%%%%%%%%%%%%%%%%%%%%%%%%%%%%%%%%%%%%%%%%%%%%%%%%%%%%%%%%%%

    \textbf{CPU} & \textbf{\hyperref[regdesc:SYSTEMCPUPERICLKENREG]{SYSTEM\_CPU\_PERI\_CLK\_EN\_REG}} & \textbf{\hyperref[regdesc:SYSTEMCPUPERIRSTENREG]{SYSTEM\_CPU\_PERI\_RST\_EN\_REG}} \\ \hline
    DEBUG\_ASSIST & \hyperref[fielddesc:SYSTEMCLKENASSISTDEBUG]{SYSTEM\_CLK\_EN\_ASSIST\_DEBUG} & \hyperref[fielddesc:SYSTEMRSTENASSISTDEBUG]{SYSTEM\_RST\_EN\_ASSIST\_DEBUG} \\\hline

%%%%%%%%%%%%%%%%%%%%%%%%%%%%%%%%%%%%%%%%%%%%%%%%%%%%%%%%%%%%%%%%%%%%%%%%%%%%%%%%%%%%%%%%%%%%%%%%%%%%%%%%%%%%%%%%%%%%%%%%%%%

%%%%%%%%%%%%%%%%%%%%%%%%%%%%%%%%%%%%%%%%%%%%%%%%%%%%%%%%%%%%%%%%%%%%%%%%%%%%%%%%%%%%%%%%%%%%%%%%%%%%%%%%%%%%%%%%%%%%%%%%%%%
    \textbf{Peripherals} & \textbf{\hyperref[regdesc:SYSTEMPERIPCLKEN0REG]{SYSTEM\_PERIP\_CLK\_EN0\_REG}} &  \textbf{\hyperref[regdesc:SYSTEMPERIPRSTEN0REG]{SYSTEM\_PERIP\_RST\_EN0\_REG}} \\ \hline

    TIMER    & \hyperref[fielddesc:SYSTEMTIMERSCLKEN]{SYSTEM\_TIMERS\_CLK\_EN} & \hyperref[fielddesc:SYSTEMTIMERSRST]{SYSTEM\_TIMERS\_RST} \\\hline
    SPI0 / SPI1     & \hyperref[fielddesc:SYSTEMSPI01CLKEN]{SYSTEM\_SPI01\_CLK\_EN}         & \hyperref[fielddesc:SYSTEMSPI01RST]{SYSTEM\_SPI01\_RST} \\\hline
    UART0           & \hyperref[fielddesc:SYSTEMUARTCLKEN]{SYSTEM\_UART\_CLK\_EN}          & \hyperref[fielddesc:SYSTEMUARTRST]{SYSTEM\_UART\_RST} \\\hline
    UART1           & \hyperref[fielddesc:SYSTEMUART1CLKEN]{SYSTEM\_UART1\_CLK\_EN}         & \hyperref[fielddesc:SYSTEMUART1RST]{SYSTEM\_UART1\_RST} \\\hline
    I2S             & \hyperref[fielddesc:SYSTEMI2S0CLKEN]{SYSTEM\_I2S0\_CLK\_EN}         & \hyperref[fielddesc:SYSTEMI2S0RST]{SYSTEM\_I2S0\_RST} \\\hline
    SPI2            & \hyperref[fielddesc:SYSTEMSPI2CLKEN]{SYSTEM\_SPI2\_CLK\_EN}          & \hyperref[fielddesc:SYSTEMSPI2RST]{SYSTEM\_SPI2\_RST} \\\hline
    I2C0            & \hyperref[fielddesc:SYSTEMEXT0CLKEN]{SYSTEM\_EXT0\_CLK\_EN}          & \hyperref[fielddesc:SYSTEMEXT0RST]{SYSTEM\_EXT0\_RST} \\\hline
    UHCI0           & \hyperref[fielddesc:SYSTEMUHCI0CLKEN]{SYSTEM\_UHCI0\_CLK\_EN}         & \hyperref[fielddesc:SYSTEMUHCI0RST]{SYSTEM\_UHCI0\_RST} \\\hline
    RMT             & \hyperref[fielddesc:SYSTEMRMTCLKEN]{SYSTEM\_RMT\_CLK\_EN}           & \hyperref[fielddesc:SYSTEMRMTRST]{SYSTEM\_RMT\_RST} \\\hline
    LED PWM Controller & \hyperref[fielddesc:SYSTEMLEDCCLKEN]{SYSTEM\_LEDC\_CLK\_EN}          & \hyperref[fielddesc:SYSTEMLEDCRST]{SYSTEM\_LEDC\_RST} \\\hline
    Timer Group0    & \hyperref[fielddesc:SYSTEMTIMERGROUPCLKEN]{SYSTEM\_TIMERGROUP\_CLK\_EN} & \hyperref[fielddesc:SYSTEMTIMERGROUPRST]{SYSTEM\_TIMERGROUP\_RST} \\\hline
    Timer Group1    & \hyperref[fielddesc:SYSTEMTIMERGROUP1CLKEN]{SYSTEM\_TIMERGROUP1\_CLK\_EN} & \hyperref[fielddesc:SYSTEMTIMERGROUP1RST]{SYSTEM\_TIMERGROUP1\_RST} \\\hline
    TWAI Controller    & \hyperref[fielddesc:SYSTEMCANCLKEN]{SYSTEM\_CAN\_CLK\_EN} & \hyperref[fielddesc:SYSTEMCANRST]{SYSTEM\_CAN\_RST} \\\hline
    USB\_DEVICE    & \hyperref[fielddesc:SYSTEMUSBDEVICECLKEN]{SYSTEM\_USB\_DEVICE\_CLK\_EN} & \hyperref[fielddesc:SYSTEMUSBDEVICERST]{SYSTEM\_USB\_DEVICE\_RST} \\\hline
    UART MEM    & \hyperref[fielddesc:SYSTEMUARTMEMCLKEN]{SYSTEM\_UART\_MEM\_CLK\_EN} \tnote{4} & \hyperref[fielddesc:SYSTEMUARTMEMRST]{SYSTEM\_UART\_MEM\_RST} \\\hline
    APB SARADC    & \hyperref[fielddesc:SYSTEMAPBSARADCCLKEN]{SYSTEM\_APB\_SARADC\_CLK\_EN} & \hyperref[fielddesc:SYSTEMAPBSARADCRST]{SYSTEM\_APB\_SARADC\_RST} \\\hline
    ADC Controller    & \hyperref[fielddesc:SYSTEMADC2ARBCLKEN]{SYSTEM\_ADC2\_ARB\_CLK\_EN} & \hyperref[fielddesc:SYSTEMADC2ARBRST]{SYSTEM\_ADC2\_ARB\_RST} \\\hline
    System Timer    & \hyperref[fielddesc:SYSTEMSYSTIMERCLKEN]{SYSTEM\_SYSTIMER\_CLK\_EN} & \hyperref[fielddesc:SYSTEMSYSTIMERRST]{SYSTEM\_SYSTIMER\_RST} \\\hline
%%%%%%%%%%%%%%%%%%%%%%%%%%%%%%%%%%%%%%%%%%%%%%%%%%%%%%%%%%%%%%%%%%%%%%%%%%%%%%%%%%%%%%%%%%%%%%%%%%%%%%%%%%%%%%%%%%%%%%%%%%%

%%%%%%%%%%%%%%%%%%%%%%%%%%%%%%%%%%%%%%%%%%%%%%%%%%%%%%%%%%%%%%%%%%%%%%%%%%%%%%%%%%%%%%%%%%%%%%%%%%%%%%%%%%%%%%%%%%%%%%%%%%%
    \textbf{Accelerators} & \textbf{\hyperref[regdesc:SYSTEMPERIPCLKEN1REG]{SYSTEM\_PERIP\_CLK\_EN1\_REG}} & \textbf{\hyperref[regdesc:SYSTEMPERIPRSTEN1REG]{SYSTEM\_PERIP\_RST\_EN1\_REG}} \\ \hline

    TSENS & \hyperref[fielddesc:SYSTEMTSENSCLKEN]{SYSTEM\_TSENS\_CLK\_EN} & \hyperref[fielddesc:SYSTEMTSENSRST]{SYSTEM\_TSENS\_RST} \\\hline
    DMA & \hyperref[fielddesc:SYSTEMDMACLKEN]{SYSTEM\_DMA\_CLK\_EN} & \hyperref[fielddesc:SYSTEMDMARST]{SYSTEM\_DMA\_RST}\tnote{5}  \\\hline
    HMAC & \hyperref[fielddesc:SYSTEMCRYPTOHMACCLKEN]{SYSTEM\_CRYPTO\_HMAC\_CLK\_EN} & \hyperref[fielddesc:SYSTEMCRYPTOHMACRST]{SYSTEM\_CRYPTO\_HMAC\_RST} \tnote{6} \\\hline
    Digital Signature & \hyperref[fielddesc:SYSTEMCRYPTODSCLKEN]{SYSTEM\_CRYPTO\_DS\_CLK\_EN} & \hyperref[fielddesc:SYSTEMCRYPTODSRST]{SYSTEM\_CRYPTO\_DS\_RST} \tnote{7} \\\hline
    RSA Accelerator & \hyperref[fielddesc:SYSTEMCRYPTORSACLKEN]{SYSTEM\_CRYPTO\_RSA\_CLK\_EN} & \hyperref[fielddesc:SYSTEMCRYPTORSARST]{SYSTEM\_CRYPTO\_RSA\_RST} \\\hline
    SHA Accelerator & \hyperref[fielddesc:SYSTEMCRYPTOSHACLKEN]{SYSTEM\_CRYPTO\_SHA\_CLK\_EN} & \hyperref[fielddesc:SYSTEMCRYPTOSHARST]{SYSTEM\_CRYPTO\_SHA\_RST} \\\hline
    AES Accelerator & \hyperref[fielddesc:SYSTEMCRYPTOAESCLKEN]{SYSTEM\_CRYPTO\_AES\_CLK\_EN} & \hyperref[fielddesc:SYSTEMCRYPTOAESRST]{SYSTEM\_CRYPTO\_AES\_RST} \\\hline

\end{longtable}
\end{ThreePartTable}


% >>> If your module has no registers, delete sample text
%       for Sections Register Summary and Registers below <<<

%%%%%%%%%%%%%%%%%%%%%%%%
%%% Register Summary %%%
%%%%%%%%%%%%%%%%%%%%%%%%

% >>> Update the sample text below and insert
%     the info on register summary into the subfile <<<

\clearpage
\section{Register Summary}
\label{sec:sysreg-reg-summ}
\hypertarget{sysreg-reg-summ}{}

The addresses in this section are relative to the base address of system registers provided in Table \ref{tab:sysmem-base-address} in Chapter \ref{mod:sysmem} \textit{\nameref{mod:sysmem}}.

The abbreviations given in Column \textbf{Access} are explained in Section \hyperref[glossary-access-types]{\textit{Access Types for Registers}}.

\subfile{\modulefiles/18-SYSREG-reg-summ__EN}

The addresses below are relative to the base address of apb control provided in Table \ref{tab:sysmem-base-address} in Chapter \ref{mod:sysmem} \textit{\nameref{mod:sysmem}}.

The abbreviations given in Column \textbf{Access} are explained in Section \hyperref[glossary-access-types]{\textit{Access Types for Registers}}.

\subfile{\modulefiles/18-SYSREG-apb-reg-summ__EN}

%%%%%%%%%%%%%%%%%
%%% Registers %%%
%%%%%%%%%%%%%%%%%

% >>> Update the sample text below and insert
%      the info on registers into the subfile <<<

\clearpage
\section{Registers}

The addresses below are relative to the base address of system register provided in Table \ref{tab:sysmem-base-address} in Chapter \ref{mod:sysmem} \textit{\nameref{mod:sysmem}}.

\subfile{\modulefiles/18-SYSREG-reg__EN}

The addresses below are relative to the base address of apb control provided in Table \ref{tab:sysmem-base-address} in Chapter \ref{mod:sysmem} \textit{\nameref{mod:sysmem}}.

\subfile{\modulefiles/18-SYSREG-apb-reg__EN}

\end{document}
